\section{Introduction}
\label{sec:intro}
% Five moves, knowledge/paper-anatomy.md Section 2. No subsections, use \paragraph{}.

% (1) Cost. A real incident or survey number, then the open problem.
<Problem> is pervasive and costly~\cite{placeholder2024}. In <year>, <concrete incident with a number>~\cite{placeholder2024}. How to <goal> when only <constraint> therefore remains an open problem.

% (2) Shift of view and the named properties.
Yet <traditional objective> does not guarantee <downstream goal>. Under this view, two properties govern <outcome>. \emph{<Property A>} measures <...>. \emph{<Property B>} measures <...>. Figure~\ref{fig:motivation} and Example~\ref{ex:motivation} make the trade-off concrete.

\begin{figure}[t]
\centering
\fbox{\parbox{0.9\linewidth}{\centering\vspace{1.2cm}Motivation figure, see knowledge/figure-archetypes.md, Section 1.\vspace{1.2cm}}}
\caption{<Under which condition what is compared, one natural sentence per panel>.}
\label{fig:motivation}
\end{figure}

% (3) Example and Observations on a real dataset.
\begin{example}[<Short title>]
\label{ex:motivation}
We illustrate with a subset of the public <Dataset>~\cite{placeholder2024}. <What each row is, the task and the model>. <Split and sizes>. <Search space and what Figure~\ref{fig:motivation} shows>.
\end{example}

\textbf{Observations.} Two observations follow from this example.
\textbf{(1) <One-sentence claim>.} <Evidence with numbers>.
\textbf{(2) <One-sentence claim>.} <Evidence with numbers>.

% (4) Challenges. One bridging question, then numbered challenges.
\paragraph{Challenges.} The observations above suggest that <requirement>. Why, then, have existing methods not addressed <problem>? We trace this to three challenges.
\textbf{(1) <Benefit is task-dependent>.} <Why, with Example~\ref{ex:motivation}>. Existing <family>~\cite{placeholder2024} lack <capability>.
\textbf{(2) <Decision space is combinatorial>.} <Why>.
\textbf{(3) <Decisions do not transfer>.} <Why>.

% (5) Contributions. The first carries the core claim, the last reports results.
\paragraph{Contributions.} To address these shortcomings, we propose \sys{}, <one-line definition>, with the following contributions.
(1) \textbf{<Name>.} <What and why, Theorem~\ref{thm:guarantee}>.
(2) \textbf{<Name>.} <What and why, Section~\ref{sec:method}>.
(3) \textbf{Results.} We evaluate \sys{} on \num{datasets} datasets against \num{baselines} baselines. \sys{} <headline result>, using <cost>.
